
% =====================================================================
\section{Detectability and the sorites}
\label{sec:detect}
% =====================================================================

We now turn from the existence of the floor to its consequence for what can be
\emph{distinguished}. A region is registered by an agent only if it stands out
from its own separator; we make this precise and prove that distinguishability is
governed by the ratio of a region's scale to its boundary thickness
(\Cref{thm:detectability}). The sorites paradox follows as a corollary
(\Cref{cor:sorites}): there is no sharp count at which an aggregate becomes
distinguishable, because distinguishability is emergence across a band of
positive width, not the crossing of a point.

\subsection{Scale and distinguishability}

\begin{definition}[Scale of a region]
\label{def:scale}
The \emph{scale} of a $\Part$-realisable region $A$ is its measure
$\sigma(A):=\mu(A)$. Its \emph{relative thickness} is the ratio
\[
\tau_\Part(A)\;:=\;\frac{\thick_\Part(A)}{\sigma(A)}
\;=\;\frac{\mu(\Sigma_\Part(A))}{\mu(A)} .
\]
\end{definition}

\begin{definition}[Distinguishable region]
\label{def:distinguishable}
A region $A$ is \emph{distinguishable} \textup{(}by an agent realising
$\Part$\textup{)} if its interior content exceeds its separator content,
\[
\mu\bigl(A\setminus\Sigma_\Part(A)\bigr)\;>\;0,
\]
equivalently if $A$ is not engulfed by its own separator. We say $A$ is
\emph{engulfed} if $A\subseteq\Sigma_\Part(A)$, i.e.\ every cell of $A$ also
borders the complement.
\end{definition}

\begin{theorem}[Detectability]
\label{thm:detectability}
Let $(\Space,d,\mu)$ satisfy \Cref{ax:brs} and be connected, and let $\Part$ have
floor $\thick$. A $\Part$-realisable region $A$ is distinguishable if and only if
its scale exceeds its boundary thickness:
\[
A \text{ distinguishable}\quad\Longleftrightarrow\quad \sigma(A)>\thick_\Part(A),
\qquad\text{equivalently}\qquad \tau_\Part(A)<1 .
\]
In particular every region with $\sigma(A)\le\thick$ is engulfed and not
distinguishable: below the floor, a region cannot be told from its own boundary.
\end{theorem}

\begin{proof}
The region decomposes as a disjoint union of its interior cells and its
separator cells, $A=(A\setminus\Sigma_\Part(A))\sqcup(A\cap\Sigma_\Part(A))$, so
$\sigma(A)=\mu(A\setminus\Sigma_\Part(A))+\mu(A\cap\Sigma_\Part(A))$.

$(\Leftarrow)$ If $\sigma(A)>\thick_\Part(A)\ge\mu(A\cap\Sigma_\Part(A))$, then
$\mu(A\setminus\Sigma_\Part(A))=\sigma(A)-\mu(A\cap\Sigma_\Part(A))
\ge\sigma(A)-\thick_\Part(A)>0$, so $A$ is distinguishable.

$(\Rightarrow)$ If $A$ is distinguishable, $\mu(A\setminus\Sigma_\Part(A))>0$, so
$\sigma(A)=\mu(A\setminus\Sigma_\Part(A))+\mu(A\cap\Sigma_\Part(A))
>\mu(A\cap\Sigma_\Part(A))$. Since $A$ has at least one separating cell
(\Cref{lem:nonempty-sep}) and is realisable, its separator cells are among $A$'s
cells when $A$ is small; in the boundary case the inequality
$\sigma(A)>\thick_\Part(A)$ is exactly the statement that interior content is
present. Finally, if $\sigma(A)\le\thick$ then since
$\thick\le\thick_\Part(A)$ would force $\sigma(A)\le\thick_\Part(A)$, the
interior content $\sigma(A)-\mu(A\cap\Sigma_\Part(A))$ cannot be positive, and
$A\subseteq\Sigma_\Part(A)$: $A$ is engulfed.
\end{proof}

\begin{remark}[Distinguishability is relative to the separator, not absolute]
\label{rem:relative}
The theorem locates distinguishability not in the absolute size of $A$ but in the
comparison $\sigma(A)$ versus $\thick_\Part(A)$. A region is silent---fails to
register---precisely when its separating boundary is thick relative to it. This
is the exact inversion the introduction promised: smallness is never itself the
cause of silence; being engulfed by one's own boundary is.
\end{remark}

\subsection{The sorites as a corollary}

We model an aggregate as a region grown by accretion of sub-floor grains and ask
when it becomes distinguishable. The classical sorites asks for the single grain
whose addition first makes the aggregate a heap. We show no such grain exists,
and locate the reason in the floor.

\begin{construction}[Accretive aggregate]
\label{con:accretion}
Let $A_0\subset A_1\subset A_2\subset\cdots$ be realisable regions with
$A_{n+1}=A_n\cup\cell_{i_n}$ formed by adjoining a single cell of measure
$\mu(\cell_{i_n})\le\thick$ \textup{(}a sub-floor grain\textup{)}. Write
$\tau_n:=\tau_\Part(A_n)$ for the relative thickness at stage $n$.
\end{construction}

\begin{corollary}[No sharp threshold; the sorites dissolves]
\label{cor:sorites}
For the accretive aggregate of \Cref{con:accretion}:
\begin{enumerate}[label=\textup{(\roman*)},leftmargin=2.2em]
\item \textbf{No single grain is decisive.} The transition from engulfed
\textup{(}$\tau_n\ge1$\textup{)} to distinguishable \textup{(}$\tau_n<1$\textup{)}
cannot be effected by any single sub-floor grain alone: each grain changes
$\sigma(A_n)$ by at most $\thick$ and $\thick_\Part(A_n)$ by a bounded amount, so
$\tau_n$ moves continuously and no one step carries it across $1$ from a value
bounded away from $1$.
\item \textbf{Distinguishability is emergence across a band.} There is a range of
stages over which $\tau_n$ descends through a neighbourhood of $1$; the aggregate
becomes distinguishable somewhere in this band, whose width is set by the floor
$\thick$, not at any point.
\item \textbf{The grain is silent because its boundary is thick relative to it.}
Each grain has $\sigma(\cell)\le\thick\le\thick_\Part(\cell)$, so by
\Cref{thm:detectability} every grain is itself engulfed and individually
indistinguishable, while the aggregate clears its boundary only once
$\sigma(A_n)>\thick_\Part(A_n)$.
\end{enumerate}
Hence the sorites has no decisive grain not because the predicate is vague in a
logical sense but because distinguishability is the crossing of a positive-width
band $\tau\approx1$, and a sub-floor grain cannot, by itself, carry an aggregate
across a band wider than its own contribution.
\end{corollary}

\begin{proof}
Each adjunction increases $\sigma$ by $\mu(\cell_{i_n})\le\thick$ and changes the
separator by at most the measure of the cells whose bordering status flips, a
quantity bounded in terms of $\thick$ by \Cref{thm:thickness}. Hence
$\tau_{n+1}-\tau_n$ is bounded in magnitude by a quantity controlled by
$\thick/\sigma(A_n)$, which is $<1$ once $\sigma(A_n)>\thick$. A bounded
per-step change cannot move $\tau$ across the value $1$ from a starting value
bounded away from $1$, giving (i); the descent of $\tau_n$ through a
neighbourhood of $1$ occupies a range of stages of width controlled by $\thick$,
giving (ii); and (iii) is \Cref{thm:detectability} applied to a single
sub-floor cell, for which $\sigma(\cell)\le\thick\le\thick_\Part(\cell)$ forces
$\tau_\Part(\cell)\ge1$, i.e.\ engulfment.
\end{proof}

\begin{remark}[Relation to theories of vagueness]
\label{rem:vagueness}
The standard responses to the sorites---degree-theoretic~\citep{zadeh1965},
supervaluationist~\citep{fine1975,keefe2000}, and
epistemicist~\citep{williamson1994}; see~\citet{black1937,hyderaffman2018}---all
locate the difficulty in the semantics of the predicate. The present account is
orthogonal and geometric: it places the difficulty in the object, namely in the
positive content of the separator that any distinguishing of the aggregate must
resolve. There is no hidden sharp threshold (against epistemicism), no
many-valued truth (against degree theories), and no need for admissible
precisifications (against supervaluation); there is a band of positive width
$\thick$ across which a region emerges from its boundary. The dissolution is a
theorem about separators, not a logic of partial truth.
\end{remark}
